\documentclass[11pt,a4paper]{article}

\usepackage[a4paper,margin=2.2cm]{geometry}
\usepackage{fontspec}
\setmainfont{Times New Roman}
\usepackage{polyglossia}
\setdefaultlanguage{vietnamese}
\setotherlanguage{english}
\usepackage{amsmath,amssymb,booktabs,array,longtable}
\usepackage{hyperref}
\usepackage{xcolor}
\usepackage{listings}
\usepackage{enumitem}

\hypersetup{
  colorlinks=true,
  linkcolor=blue,
  urlcolor=blue,
  citecolor=blue
}

\lstdefinestyle{mystyle}{
  basicstyle=\ttfamily\small,
  breaklines=true,
  frame=single,
  columns=fullflexible,
  keepspaces=true,
  showstringspaces=false
}
\lstset{style=mystyle}

\title{\textbf{Ứng dụng Answer Set Programming trong 3 môn học CSE}\\
\large PPL Honors Module Mini-Project}
\author{
  \textbf{Nhóm 8} \\
  2310543 \quad HỒ ANH DŨNG \textless dung.hokhmt2k23@hcmut.edu.vn\textgreater \\
  2312264 \quad LÊ THÀNH NGHĨA \textless nghia.lethanh58566@hcmut.edu.vn\textgreater \\
  2312291 \quad VŨ TRỌNG NGHĨA \textless nghia.vutrong@hcmut.edu.vn\textgreater
}
\date{Tháng 04, 2026}

\begin{document}
\maketitle

\begin{abstract}
Trong mini-project này, nhóm em hiện thực 3 ứng dụng ASP tương ứng với 3 môn: Principles of Programming Languages (PPL), Database Systems, và Operating Systems. Với mỗi môn, nhóm em mô tả bài toán, xây dựng mô hình ASP, chạy Clingo để lấy stable model, và phân tích ý nghĩa kết quả. Phiên bản hiện tại không chỉ dừng ở ví dụ tối thiểu mà đã được mở rộng thành bộ kiểm thử gồm 10 case cho PPL, 4 case cho Database, và 4 case cho Operating Systems. Mục tiêu của nhóm là làm rõ cách tư duy khai báo trong ASP có thể áp dụng xuyên suốt nhiều miền kiến thức khác nhau, đồng thời giữ được tính kiểm chứng tự động, khả năng giải thích kết quả, và sự nhất quán trong cấu trúc thực nghiệm.
\end{abstract}

\section{Introduction}
Theo đề bài trong \texttt{main\_en.txt}, báo cáo cần có 3 ứng dụng thuộc 3 môn khác nhau (trong đó bắt buộc có PPL), đồng thời có phần so sánh liên môn. Nhóm em chọn:
\begin{itemize}[leftmargin=1.5em]
  \item PPL: phân tích \textit{scope} và \textit{binding} dưới góc nhìn ngữ nghĩa tĩnh.
  \item Database Systems: kiểm tra \textit{conflict-serializability} của lịch giao dịch.
  \item Operating Systems: kiểm tra trạng thái an toàn theo phong cách Banker's algorithm.
\end{itemize}

Điểm chung của cả 3 bài là: mô tả tri thức bằng facts/rules/constraints, sau đó dùng stable model để trả lời câu hỏi đúng/sai hoặc sinh một phương án hợp lệ. Điểm khác nhau nằm ở loại đối tượng được mô hình hóa: quan hệ ngữ nghĩa chương trình, đồ thị phụ thuộc giao dịch, và động học tài nguyên trong hệ điều hành.

\section{Background}
\subsection{ASP và stable model}
Answer Set Programming (ASP) là một mô hình lập trình khai báo dựa trên stable model semantics \cite{gelfond1988,brewka2011}. Thay vì viết thuật toán theo từng bước như phong cách mệnh lệnh, ta mô tả bài toán bằng các quan hệ logic; solver (ở đây là Clingo \cite{gebser2011}) sẽ tìm các mô hình thỏa mãn toàn bộ ràng buộc.

Đối với mini-project này, ASP phù hợp vì cả ba miền đều có thể được phát biểu bằng quan hệ và ràng buộc:
\begin{itemize}[leftmargin=1.5em]
  \item PPL yêu cầu suy ra thuộc tính ngữ nghĩa tĩnh của tên định danh.
  \item Database yêu cầu suy diễn cạnh phụ thuộc và phát hiện chu trình.
  \item Operating Systems yêu cầu kiểm tra sự tồn tại của một thứ tự chạy thỏa ràng buộc tài nguyên.
\end{itemize}

\subsection{Tổ chức mã nguồn và thực nghiệm}
Để làm submission chuyên nghiệp hơn, nhóm em tổ chức lại mã nguồn thành:
\begin{itemize}[leftmargin=1.5em]
  \item \texttt{scripts/}: các Python runner cho 3 ứng dụng.
  \item \texttt{asp/}: các file ASP baseline được sinh ra.
  \item \texttt{cases/}: các case mở rộng, phân theo từng môn.
  \item \texttt{report/}: mã nguồn và PDF của báo cáo.
\end{itemize}
Cách tổ chức này giúp tách biệt rõ mã điều phối, mã ASP sinh ra, và tài liệu báo cáo.

\section{Main Content: Three ASP Applications}
\subsection{Môn 1: PPL -- Scope and Binding Analysis}
\subsubsection*{Bài toán}
Cho tập biến, khai báo biến trong hàm, và kiểu dữ liệu, nhóm em cần suy ra:
\begin{itemize}[leftmargin=1.5em]
  \item Biến nào là \texttt{local\_variable}, biến nào là \texttt{global\_variable}.
  \item Biến thuộc scope nào (\texttt{in\_scope}).
  \item Biến nào được bind hợp lệ với kiểu \texttt{int} (\texttt{bound\_variable}).
\end{itemize}
Ngoài ra nhóm em thêm các predicate chẩn đoán lỗi thường gặp trong bảng ký hiệu trung gian: \texttt{missing\_scope}, \texttt{untyped\_variable}, \texttt{orphan\_declaration}, và \texttt{non\_int\_type}.

Để tăng \textit{PPL depth} so với mức mô hình hóa theo hàm (function-level) đơn giản, nhóm em mở rộng mô hình theo hướng \textit{lexical scoping}: hỗ trợ scope lồng nhau (nested scopes), \textit{shadowing} (khai báo bên trong che khuất bên ngoài), và phát hiện \textit{free variable} (biến được dùng nhưng không có khai báo nào truy cập được).

\subsubsection*{Mô hình ASP cốt lõi}
\begin{lstlisting}[language={}]
global_variable(X) :- variable(X), not local_variable(X, _).
local_variable(X, F) :- declared_in(X, F).

in_scope(X, S) :- local_variable(X, F), function_scope(F, S).
in_scope(X, global) :- global_variable(X).

bound_variable(X, S) :- in_scope(X, S), type_of(X, int).

orphan_declaration(X, F) :- declared_in(X, F), not variable(X).
missing_scope(F) :- declared_in(_, F), not function_scope(F, _).
untyped_variable(X) :- variable(X), not type_of(X, _).
non_int_type(X, T) :- type_of(X, T), T != int.
\end{lstlisting}

\subsubsection*{Mở rộng: lexical scope, shadowing, free variable}
Các fact bổ sung (tùy chọn) để mô hình hóa lexical scoping:
\begin{itemize}[leftmargin=1.5em]
  \item \texttt{scope\_parent(Inner, Outer)}: quan hệ bao-hàm giữa hai scope.
  \item \texttt{declared\_in\_scope(X,S)}: biến \texttt{X} được khai báo tại scope \texttt{S}.
  \item \texttt{use\_in(X,S)}: biến \texttt{X} được sử dụng tại scope \texttt{S}.
\end{itemize}

Ý tưởng chính là: một \texttt{use\_in(X,S)} sẽ bind tới \textit{khai báo gần nhất} theo chuỗi \texttt{scope\_parent} (closest enclosing declaration). Nếu không có khai báo nào truy cập được thì \texttt{X} là free variable tại \texttt{S}.

\begin{lstlisting}[language={}]
ancestor_or_self(S, S) :- scope(S).
ancestor_or_self(S, A) :- scope_parent(S, P), ancestor_or_self(P, A).

reachable_decl(UseS, X, DeclS) :- ancestor_or_self(UseS, DeclS), declared_in_scope(X, DeclS).
shadowed_by(UseS, X, DeclS) :- reachable_decl(UseS, X, DeclS), reachable_decl(UseS, X, DeclS2),
                              ancestor_or_self(DeclS2, DeclS), DeclS2 != DeclS.
binding(UseS, X, DeclS) :- reachable_decl(UseS, X, DeclS), not shadowed_by(UseS, X, DeclS).

resolved_use(X, UseS, DeclS) :- use_in(X, UseS), binding(UseS, X, DeclS).
free_var(X, UseS) :- use_in(X, UseS), not binding(UseS, X, _).
\end{lstlisting}

\subsubsection*{Ý nghĩa PPL}
Về mặt lý thuyết ngôn ngữ, mô hình này đang xấp xỉ một bài toán phân tích ngữ nghĩa tĩnh trên bảng ký hiệu:\
\begin{itemize}[leftmargin=1.5em]
  \item \texttt{declared\_in/2} biểu diễn nơi khai báo của tên định danh.
  \item \texttt{function\_scope/2} ánh xạ hàm sang scope trừu tượng mà solver có thể reason.
  \item \texttt{bound\_variable/2} chỉ được suy ra nếu thông tin scope và type đều nhất quán.
\end{itemize}
Điểm quan trọng là mô hình đang giả định \textit{static binding} và làm việc ở mức khai báo-hàm, chứ chưa mô tả đầy đủ block nesting, lambda hoặc cơ chế lexical environment chi tiết. Tuy nhiên, trong khuôn khổ bài này, nó vẫn thể hiện đúng tinh thần PPL: dùng luật khai báo để kiểm chứng thuộc tính ngữ nghĩa thay vì duyệt cây cú pháp bằng thuật toán mệnh lệnh.

\subsubsection*{Kết quả thực nghiệm}
 Phiên bản hiện tại chạy \texttt{12/12} case PASS bằng script \texttt{scripts/scope\_binding.py}. Ngoài case baseline, nhóm em dùng thêm các tình huống đáng chú ý:
\begin{itemize}[leftmargin=1.5em]
  \item \textbf{global\_only}: biến không có khai báo hàm phải rơi về global scope.
  \item \textbf{missing\_function\_scope}: phát hiện thông tin bảng ký hiệu bị thiếu ánh xạ scope.
  \item \textbf{multi\_decl}: cùng một biến xuất hiện ở nhiều khai báo hàm, cho thấy mô hình hiện tại có thể biểu diễn sự nhập nhằng hoặc tái sử dụng tên.
  \item \textbf{global\_non\_int}: biến toàn cục vẫn có scope hợp lệ nhưng không được bind nếu kiểu khác \texttt{int}.
  \item \textbf{mixed\_diagnostics}: một case tổng hợp vừa có local binding hợp lệ, vừa có global non-int, vừa có orphan declaration.
  \item \textbf{shadowing\_nested\_scope}: kiểm tra shadowing giữa scope lồng nhau và ràng buộc “closest declaration wins”.
  \item \textbf{free\_variable\_use}: phát hiện free variable khi có \texttt{use\_in} nhưng không có \texttt{declared\_in\_scope} phù hợp.
\end{itemize}

Ví dụ stable model của case tổng hợp:
\begin{lstlisting}[language={}]
global_variable(b) bound_variable(a,f1) in_scope(a,f1)
in_scope(c,f1) in_scope(b,global) local_variable(a,foo)
local_variable(c,foo) non_int_type(b,string)
orphan_declaration(c,foo)
\end{lstlisting}

Stable model trên cho thấy ba loại thông tin cùng tồn tại trong một lời giải: (1) binding đúng của \texttt{a}; (2) biến \texttt{b} có scope nhưng sai kiểu nên không được bind; (3) \texttt{c} tạo ra cảnh báo orphan declaration do có khai báo nhưng thiếu fact \texttt{variable(c)}. Điều này quan trọng hơn ví dụ baseline vì nó mô phỏng một trạng thái phân tích bán-lỗi trong frontend compiler.

\subsubsection*{Hạn chế}
Mô hình PPL hiện tại vẫn còn đơn giản ở ba điểm:
\begin{itemize}[leftmargin=1.5em]
  \item Lexical scoping hiện được mô hình hóa ở mức quan hệ \texttt{scope\_parent} trừu tượng, chưa gắn trực tiếp với cây cú pháp (AST) và block nesting sinh ra từ cú pháp thật.
  \item Free-variable hiện được phát hiện từ fact \texttt{use\_in/2} (mức bảng ký hiệu/IR), chưa mô hình hóa dùng biến bên trong biểu thức cụ thể.
  \item Predicate \texttt{multi\_decl} mới phản ánh khai báo lặp, chưa đủ để mô hình hóa capture hay alpha-renaming.
\end{itemize}
Tuy vậy, trong khuôn khổ mini-project, phần này đã có chiều sâu hơn mức ví dụ đồ chơi vì nó xử lý được trường hợp dữ liệu ngữ nghĩa thiếu, sai kiểu, hoặc không nhất quán.

\subsection{Môn 2: Database Systems -- Conflict-Serializability}
\subsubsection*{Bài toán}
Với một schedule gồm các thao tác read/write trên nhiều transaction, nhóm em cần:
\begin{itemize}[leftmargin=1.5em]
  \item Xây precedence graph từ xung đột dữ liệu.
  \item Kiểm tra có chu trình hay không.
  \item Nếu không chu trình thì sinh một serial order hợp lệ.
\end{itemize}

\subsubsection*{Mô hình ASP cốt lõi}
\begin{lstlisting}[language={}]
conflict(I,J,T1,T2,X) :- op(I,T1,w,X), op(J,T2,_,X), T1 != T2, I < J.
conflict(I,J,T1,T2,X) :- op(I,T1,_,X), op(J,T2,w,X), T1 != T2, I < J.
edge(T1,T2) :- conflict(_,_,T1,T2,_).

path(T1,T2) :- edge(T1,T2).
path(T1,T3) :- path(T1,T2), edge(T2,T3).
cycle(T) :- path(T,T).

serializable :- not cycle(_).
not_serializable :- cycle(_).
\end{lstlisting}

Nếu schedule acyclic, mô hình còn dùng một lớp guess-and-check để sinh một thứ tự \texttt{order(T,P)} sao cho mọi cạnh \texttt{edge(T1,T2)} đều được tôn trọng. Về bản chất, đây là bài toán topo-sort trên precedence graph được phát biểu theo dạng ràng buộc logic.

\subsubsection*{Mở rộng: recoverable/cascadeless/strict (dựa trên commit)}
Để tránh việc bài toán Database bị quá đơn giản (chỉ dừng ở conflict-serializability), nhóm em mở rộng thêm phân tích các tính chất lịch thường gặp trong DBMS dựa trên thao tác \texttt{commit}. Input được phép có thêm fact \texttt{op(I,T,c,none)}.

Ý tưởng là suy ra quan hệ \texttt{rf/5} (read-from) theo kiểu “đọc giá trị từ lần ghi gần nhất trước đó”:
\begin{lstlisting}[language={}]
rf(Tr,Tw,X,Iw,Ir) :- write_pos(Tw,X,Iw), read_pos(Tr,X,Ir), Tw != Tr, Iw < Ir,
                     not later_write_between(Iw,Ir,X).
\end{lstlisting}

Từ \texttt{rf}, ta kiểm tra:
\begin{itemize}[leftmargin=1.5em]
  \item \textbf{Recoverable}: nếu \texttt{Tr} đọc từ \texttt{Tw} và \texttt{Tr} commit, thì \texttt{Tw} phải commit trước \texttt{Tr}.
  \item \textbf{Cascadeless}: mọi đọc-from chỉ hợp lệ nếu writer đã commit trước thời điểm đọc (tránh dirty read).
  \item \textbf{Strict}: sau khi \texttt{Tw} write(X), không tx khác được read/write X trước khi \texttt{Tw} commit (mạnh hơn cascadeless).
\end{itemize}

Việc thêm 3 tính chất này giúp phần Database có chiều sâu hơn: cùng một schedule có thể conflict-serializable nhưng vẫn \textit{không} recoverable/cascadeless/strict.

\subsubsection*{Kết quả thực nghiệm}
 Phiên bản hiện tại chạy \texttt{6/6} case PASS bằng script \texttt{scripts/db\_serializability.py}. Bộ case bao gồm:
\begin{itemize}[leftmargin=1.5em]
  \item \textbf{serializable\_chain}: chuỗi phụ thuộc đơn giản \(\texttt{t1} \rightarrow \texttt{t2} \rightarrow \texttt{t3}\).
  \item \textbf{serializable\_branching}: 4 transaction với phụ thuộc phân nhánh nhưng không có chu trình.
  \item \textbf{non\_serializable\_cycle}: chu trình trực tiếp giữa 2 transaction.
  \item \textbf{non\_serializable\_three\_way}: chu trình dài \(\texttt{t1} \rightarrow \texttt{t2} \rightarrow \texttt{t3} \rightarrow \texttt{t1}\).
  \item \textbf{recoverable\_violation}: lịch vẫn serializable nhưng vi phạm recoverable/cascadeless/strict do dirty read + commit sớm.
  \item \textbf{strict\_violation}: 2 writer chồng lấp trên cùng item trước commit của writer đầu, nên không strict.
\end{itemize}

Case \textbf{serializable\_branching} sinh ra stable model sau:
\begin{lstlisting}[language={}]
serializable edge(t1,t2) edge(t1,t3) edge(t2,t4) edge(t3,t4)
order(t1,1) order(t3,2) order(t2,3) order(t4,4)
\end{lstlisting}

Kết quả này cho thấy ASP không chỉ xác định lịch là serializable mà còn tự chọn một linearization hợp lệ. Điểm đáng chú ý là \texttt{t2} và \texttt{t3} không có cạnh ràng buộc trực tiếp nên solver có thể đặt \texttt{t3} trước \texttt{t2}; điều đó phản ánh đúng thực tế rằng serial order không nhất thiết là duy nhất.

Trong case \textbf{non\_serializable\_three\_way}, stable model có dạng:
\begin{lstlisting}[language={}]
not_serializable cycle(t1) cycle(t2) cycle(t3)
edge(t1,t2) edge(t2,t3) edge(t3,t1)
\end{lstlisting}
Mô hình này hữu ích hơn một kết luận đúng/sai đơn thuần vì nó chỉ ra rõ cấu trúc gây lỗi nằm ở đâu: cả ba transaction đều thuộc một strongly connected dependency loop. Điều này minh họa lợi thế giải thích được của ASP trong các bài toán kiểm chứng.

\subsubsection*{Hạn chế}
Mô hình Database hiện tại mới dừng ở conflict-serializability:
\begin{itemize}[leftmargin=1.5em]
  \item Chưa mô hình hóa abort và các khái niệm liên quan (ví dụ strictness dưới abort, cascading abort thực sự).
  \item Chưa benchmark trên lịch lớn để đánh giá chi phí suy luận khi closure tăng.
\end{itemize}
Tuy nhiên, phần precedence graph và topo-order đã đủ vượt mức toy example vì mô hình xử lý được cả trường hợp branching và chu trình dài hơn 2 nút.

\subsection{Môn 3: Operating Systems -- Safe Sequence (Banker style)}
\subsubsection*{Bài toán}
Cho các tiến trình, tài nguyên, ma trận \texttt{alloc}, \texttt{max}, và tài nguyên hiện có \texttt{avail0}, nhóm em cần kiểm tra có tồn tại thứ tự chạy an toàn hay không.

\subsubsection*{Mô hình ASP cốt lõi}
\begin{lstlisting}[language={}]
need(P,R,N) :- max(P,R,M), alloc(P,R,A), N = M - A.

1 { run(S,P) : proc(P) } 1 :- step(S).
1 { run(S,P) : step(S) } 1 :- proc(P).

finished_before(S,P) :- step(S), run(S0,P), S0 < S.
returned(S,R,Ret) :- step(S), res(R), Ret = #sum { A,P : alloc(P,R,A),
                                  finished_before(S,P) }.
available(S,R,V) :- step(S), res(R), avail0(R,B), returned(S,R,Ret), V = B + Ret.

:- run(S,P), need(P,R,N), available(S,R,V), N > V.
\end{lstlisting}

Mô hình này thể hiện đúng bản chất guess-and-check của Banker: solver đoán một hoán vị tiến trình, sau đó loại bỏ mọi hoán vị mà tại một bước nào đó tồn tại nhu cầu tài nguyên vượt quá mức khả dụng.

\subsubsection*{Mở rộng: chẩn đoán nguyên nhân UNSAT (blocker ban đầu)}
Một hạn chế thực tế của việc chỉ in \texttt{UNSATISFIABLE} là người đọc không biết \textit{vì sao} trạng thái unsafe. Nhóm em bổ sung các predicate chẩn đoán ở bước đầu:
\begin{itemize}[leftmargin=1.5em]
  \item \texttt{lacks\_resource\_now(P,R)}: tại trạng thái ban đầu, tiến trình \texttt{P} thiếu tài nguyên \texttt{R} (need > avail0).
  \item \texttt{can\_run\_now(P)} và \texttt{blocked\_initial(P)}: phân loại tiến trình có thể/không thể là bước chạy đầu tiên.
\end{itemize}
Các atom này không thay đổi điều kiện an toàn, nhưng giúp giải thích nhanh các case unsafe: nếu \textit{không có} \texttt{can\_run\_now/1} nào thì hệ thống bị kẹt ngay từ đầu.

\subsubsection*{Kết quả thực nghiệm}
Phiên bản hiện tại chạy \texttt{4/4} case PASS bằng script \texttt{scripts/os\_banker.py}. Bộ case gồm:
\begin{itemize}[leftmargin=1.5em]
  \item \textbf{safe\_state}: ví dụ nhỏ 3 tiến trình, 2 loại tài nguyên.
  \item \textbf{safe\_state\_classic}: ví dụ lớn hơn kiểu textbook với 5 tiến trình, 3 loại tài nguyên.
  \item \textbf{unsafe\_state}: trạng thái không có tiến trình nào mở đầu được.
  \item \textbf{unsafe\_after\_partial\_progress}: có một tiến trình chạy được trước, nhưng toàn cục vẫn không an toàn.
\end{itemize}

Case \textbf{safe\_state\_classic} sinh ra stable model với chuỗi chạy an toàn:
\begin{lstlisting}[language={}]
run(1,p1) run(2,p4) run(3,p3) run(4,p0) run(5,p2)
need(p1,a,1) need(p1,b,2) need(p1,c,2)
need(p3,a,0) need(p3,b,1) need(p3,c,1)
...
\end{lstlisting}

Ý nghĩa của case này là: mặc dù có nhiều tiến trình và ba chiều tài nguyên, solver vẫn suy ra được một safe sequence hợp lệ. Đây là bước tiến đáng kể so với case nhỏ ban đầu vì không gian tìm kiếm đã tăng từ \(3!\) lên \(5!\), đồng thời predicate \texttt{available/3} phải cộng dồn tài nguyên được hoàn trả theo thời gian.

Case \textbf{unsafe\_after\_partial\_progress} đặc biệt hữu ích cho phân tích kết quả. Nó cho thấy việc tồn tại một tiến trình mở đầu hợp lệ không kéo theo trạng thái là safe. Nói cách khác, local feasibility ở bước đầu chưa đủ để kết luận global safety. Đây chính là lý do cần mô hình hóa toàn bộ thứ tự chạy bằng ASP thay vì chỉ kiểm tra greedily một bước.

\subsubsection*{Hạn chế}
Mô hình OS hiện tại vẫn còn một số đơn giản hóa:
\begin{itemize}[leftmargin=1.5em]
  \item Chưa xét tiến trình đến động theo thời gian.
  \item Chưa mô hình hóa preemption, priority, hoặc deadline.
  \item Chưa so sánh với cài đặt Banker theo phong cách thủ tục để benchmark hiệu năng.
\end{itemize}
Tuy nhiên, phần hiện tại đã đủ để thể hiện ASP như một công cụ tìm kiếm lời giải có giải thích được trong bài toán an toàn tài nguyên.

\section{Experimental Summary}
Bảng \ref{tab:coverage} tóm tắt độ phủ case hiện tại.

\begin{table}[h]
\centering
\begin{tabular}{p{2.8cm} p{2.1cm} p{5.0cm} p{3.6cm}}
\toprule
Môn & Số case & Hiện tượng chính & Mức kết quả \\
\midrule
 PPL & 12 & local/global binding, thiếu type, thiếu scope, orphan declaration, mixed diagnostics, shadowing, free variable & Stable model với atom ngữ nghĩa và atom chẩn đoán \\
 Database & 6 & precedence chain/branching/cycle và phân tích recoverable/cascadeless/strict qua commit & Serializable + serial order, kèm chẩn đoán vi phạm isolation \\
 Operating Systems & 4 & safe sequence nhỏ, safe sequence cổ điển, deadlock tức thời, deadlock sau tiến triển cục bộ & Có stable model (safe) hoặc UNSAT (unsafe) \\
\bottomrule
\end{tabular}
\caption{Độ phủ thực nghiệm của ba ứng dụng ASP}
\label{tab:coverage}
\end{table}

Bộ case trên cho thấy nhóm em đã vượt qua mức kiểm chứng tối thiểu. Thay vì chỉ chứng minh rằng chương trình chạy được trên một ví dụ, nhóm em cố gắng kiểm tra cả case đúng, case sai, case biên, và case tổng hợp. Điều này quan trọng trong rubric vì kết quả phân tích không còn là việc đọc một stable model duy nhất, mà là so sánh hành vi của mô hình trước nhiều kiểu đầu vào.

\section{Cross-domain Analysis}
\begin{table}[h]
\centering
\begin{tabular}{p{2.5cm} p{3.2cm} p{3.8cm} p{4.0cm}}
\toprule
Môn & Đối tượng facts chính & Dạng suy luận chính & Insight mô hình hóa \\
\midrule
PPL & tên định danh, khai báo, kiểu, scope & suy diễn thuộc tính ngữ nghĩa tĩnh & mạnh ở việc phát hiện dữ liệu ngữ nghĩa thiếu hoặc không nhất quán \\
Database & thao tác read/write theo thời gian & closure trên đồ thị phụ thuộc + phát hiện chu trình & tự nhiên cho bài toán đúng/sai dựa trên reachability \\
Operating Systems & cấp phát, nhu cầu cực đại, tài nguyên khả dụng & guess-and-check trên hoán vị + cộng dồn tài nguyên hoàn trả & phù hợp cho bài toán tồn tại một kế hoạch khả thi \\
\bottomrule
\end{tabular}
\caption{So sánh đặc trưng mô hình ASP giữa 3 môn}
\end{table}

Từ ba bài toán, nhóm em rút ra bốn mẫu chung khi dùng ASP:
\begin{itemize}[leftmargin=1.5em]
  \item \textbf{Tách facts và rules rõ ràng}: cả ba miền đều bắt đầu từ facts rất cụ thể, sau đó mới nâng dần thành predicate có nghĩa học thuật hơn.
  \item \textbf{Recursion là công cụ trung tâm}: ở Database là \texttt{path/2}, ở OS là tích lũy \texttt{available/3}, còn ở PPL là chuỗi suy diễn từ declaration sang scope rồi sang binding.
  \item \textbf{Constraint quyết định chất lượng mô hình}: logic lõi thường ngắn, nhưng sự đúng đắn nằm ở các constraint loại bỏ mô hình sai.
  \item \textbf{Khả năng giải thích}: output của ASP không chỉ nói \textit{đúng/sai}; nó cung cấp các atom trung gian giúp phân tích nguyên nhân của kết quả.
\end{itemize}

Sự khác nhau lớn nhất giữa ba miền nằm ở kiểu kết quả mong muốn. PPL chủ yếu suy diễn nhãn ngữ nghĩa; Database vừa kiểm tra tính chất vừa sinh topo-order; OS là bài toán tồn tại lời giải trên không gian hoán vị. Chính sự khác biệt này cho thấy ASP không bị khóa vào một mẫu duy nhất mà có thể thích ứng với nhiều kiểu bài toán trong CSE.

\section{Conclusion}
Nhóm em đã hoàn thành 3 ứng dụng ASP cho 3 môn khác nhau theo đúng yêu cầu đề bài. So với phiên bản ban đầu, dự án hiện tại có hai cải thiện quan trọng: (1) mã nguồn được tổ chức lại rõ ràng hơn cho submission; (2) bộ thực nghiệm được mở rộng đủ để cho thấy hành vi của mô hình trong nhiều tình huống, không chỉ một ví dụ minh họa.

Xét theo rubric, phần mạnh nhất của nhóm em là \textit{ASP modeling} và \textit{cross-domain consistency}. Phần cần tiếp tục cải thiện nếu muốn đẩy điểm cao hơn nữa là: mô hình PPL sâu hơn ở lexical scoping thực sự, benchmark lớn hơn cho Database/OS, và phân tích định lượng thời gian solve. Dù vậy, trong khuôn khổ mini-project hiện tại, nhóm em cho rằng dự án đã đủ tính kỹ thuật, đủ khác biệt giữa các miền, và đủ mức độ không-tầm-thường để đáp ứng tốt yêu cầu học phần.

Trong hướng mở rộng, nhóm em muốn:
\begin{itemize}[leftmargin=1.5em]
  \item thêm free-variable analysis hoặc shadowing analysis chi tiết hơn cho PPL,
  \item thêm thuộc tính recoverability và commit/abort cho Database,
  \item thêm benchmark và so sánh với phiên bản thuật toán thủ tục cho Banker,
  \item tích hợp pipeline sinh test tự động để đánh giá độ bền của mô hình.
\end{itemize}

\begin{thebibliography}{99}

\bibitem{brewka2011}
G. Brewka, T. Eiter, and M. Truszczynski.
\newblock Answer set programming at a glance.
\newblock \emph{Communications of the ACM}, 54(12):92--103, 2011.

\bibitem{gelfond1988}
M. Gelfond and V. Lifschitz.
\newblock The stable model semantics for logic programming.
\newblock In \emph{ICLP}, pages 1070--1080, 1988.

\bibitem{gebser2011}
M. Gebser et al.
\newblock Potassco: The Potsdam answer set solving collection.
\newblock \emph{AI Communications}, 24(2):107--124, 2011.

\bibitem{gabbrielli2023}
M. Gabbrielli, S. Martini, and S. Giallorenzo.
\newblock \emph{Programming Languages: Principles and Paradigms} (2nd ed.).
\newblock Springer, 2023.

\bibitem{erdem2016}
E. Erdem, M. Gelfond, and N. Leone.
\newblock Applications of answer set programming.
\newblock \emph{AI Magazine}, 37(3):53--68, 2016.

\bibitem{hahn2024}
S. Hahn et al.
\newblock Reasoning about study regulations in answer set programming.
\newblock \emph{Theory and Practice of Logic Programming}, 24(4):790--804, 2024.

\end{thebibliography}

\end{document}

